% =====================================================================
%  O texto da carta.
%
%  No padrao oficio, os paragrafos do corpo sao numerados quando o
%  documento tem mais de um — o que facilita a referencia em respostas
%  ("conforme o item 3 do oficio anterior"). Em cartas curtas ou
%  informais, a numeracao e dispensavel: basta apagar os \item e usar
%  paragrafos comuns.
% =====================================================================

% --- Vocativo --------------------------------------------------------
\noindent \lkSalutation

\vspace{1em}

% --- Corpo, com paragrafos numerados ---------------------------------
\begin{enumerate}
  \item Escreva aqui o primeiro paragrafo, apresentando o motivo do
  documento de forma direta. Quem le deve entender, ja nas primeiras
  linhas, o que se pede ou o que se comunica.

  \item Use os paragrafos seguintes para desenvolver o assunto: os fatos,
  os fundamentos e o que se espera do destinatario. Um assunto por
  paragrafo.

  \item Encerre indicando o encaminhamento desejado, o prazo, se houver, e
  o canal para eventual resposta.
\end{enumerate}

% --- Versao sem numeracao --------------------------------------------
% Para uma carta comum, apague o bloco acima e escreva assim:
%
% Escreva aqui o primeiro paragrafo.
%
% E aqui o segundo.
